%قبل از شروع مسئله ترجیحاً چیزی ننویسید. اما اگر خواستید بنویسید، در قالب لیستی (فقرات و شمارش) بنویسید.
%برای اضافه کردن نمره، بعدی از شروع‌مسئله، [x] را اضافه کنید و به جای x  نمره مورد نظر را قرار دهید. برای زیرمسئله‌ها هم به همین صورت است.
%\شروع‌مسئله[10]
\شروع‌مسئله [18]
%\متن‌سیاه{سوال عملی اول:}\\
%\متن‌سیاه{دراپ‌تست}\\
فرض کنید که \متن‌لاتین{n} ماهی داریم. ماهی \متن‌لاتین{i}-ام وزنی برابر 
$S_i$
و قدرت حمله‌ای برابر
$E_i$
دارد به صورتی که می‌دانیم: 
$E_i \le S_i$.
ماهی \متن‌لاتین{i}-ام به ماهی \متن‌لاتین{j}-ام زور می‌گوید اگر که داشته باشیم‌:
$E_i > S_j$.
حال می‌خواهیم بیشترین تعداد ماهی را انتخاب کنیم به صورتی که بین ماهی های انتخاب شده، هیچ یک از ماهی ها به ماهی دیگری زور نگوید. تعداد این ماهی ها را در 
$O(n\log(n))$
بیابید.

\solution{
پاسخ. برای حل این سوال، فرض می‌کنیم ماکسیمم قدرت حمله برابر 
$E$
باشد. بنابراین تمام ماهی‌هایی که در مجموعه هستند وزنی بیشتر یا مساوی 
$E$
دارند. زیرا در غیر این‌صورت، ماهی با قدرت حمله ماکسیمم به آنها زور می‌گوید. هم چنین اگر چنین شرطی برقرار باشد، هیچ دو ماهی‌ای به هم زور نمی‌گویند. پس کافی است ماکسیمم قدرت حمله را مشخص کنیم، سپس از بین ماهی هایی که قدرت حمله‌ای کمتر یا مساوی قدرت حمله‌ی ماکسیمم دارند، تمام ماهی هایی که وزنی بیشتر یا مساوی قدرت حمله‌ی ماکسیمم را دارند انتخاب کنیم. برای این کار ماهی ها را بر حسب قدرت حمله‌شان مرتب می‌کنیم و الگوریتم زیر را اجرا می‌کنیم:

فرض کنید ماهی‌هایی که بر حسب قدرت حمله‌شان مرتب شده‌اند،‌ شماره های ۱ تا n را به ترتیب دارند(۱ کمترین زور، n بیشترین زور). 
به ازای 
$i = 1, 2, ..., n$
مراحل زیر را اجرا می‌کنیم.

- وزن ماهی i-ام را به یک درخت AVL اضافه می‌کنیم.
 
- تعداد اعداد درخت AVL را که از قدرت حمله‌ی ماهی i-ام بیشتر است را بدست می‌آوریم.

از بین n عددی که بدست آورده ایم، بیشترین عدد، بیشترین تعداد ماهی‌ای است که می‌توانیم دور هم جمع کنیم به صورتی که هیچ دوتایی به هم زور نگویند. 

برای شمردن تعداد اعداد بزرگتر یا مساوی در درخت AVL، می‌دانیم می‌توانیم برای هر راس، تعداد اعضای زیر درخت هایش را نگه داریم. سپس اگر بخواهیم اعداد بزرگتر یا مساوی 
$x$
را حساب کنیم، x را در درخت جستجو می‌کنیم. هر بار که به بچه‌ی چپی برویم(کوچکتر یا مساوی راسی که در آن هستیم)، میفهمیم آن راس و زیر درخت سمت راست همگی بزرگتر یا مساوی هستند و آنها را با مقداری که تاکنون بدست آورده‌ایم جمع می‌زنیم. اگر هم به بچه‌ی راستی(بزرگتر از راسی که در آن هستیم) برویم چیزی به مقدارمان اضافه نمی‌کنیم. 
پس محاسبه‌ی این تعداد و هم چنین اضافه کردن اعداد به درخت AVL از 
$O(\log(n))$
است. پس کل راه از 
$O(n\log(n))$
است.

}
%\متن‌سیاه{ورودی:}\\
%در خط اول ورودی، ابتدا k و سپس m داده می‌شود.\\
%\متن‌سیاه{خروجی:}\\
%در تنها سطر خروجی، تعداد آیفون‌های مورد نیاز می‌آید. اگر در هیچ صورتی چنین کاری ممکن نبود ۱- چاپ می‌شود.\\
%\متن‌سیاه{محدودیت‌ها:}\\
%$1\leq k\leq 10^4$

%$m\leq10^9$
%\\
%\textbf{\متن‌سیاه{ورودی و خروجی نمونه}\\
%% تست اول
%\begin{latin}
%	\begin{tabularx}{\textwidth}{|X | X|}
%		\hline
%		\textbf{stdin} &
%		\textbf{stdout}\\
%		\hline
%		100 100&
%		1\\
%		\hline
%	
%	\end{tabularx}
%\end{latin}
%
%% تست دوم
%
%\begin{latin}
%	\begin{tabularx}{\textwidth}{|X | X|}
%		\hline
%		\textbf{stdin} &
%		\textbf{stdout}\\
%		\hline
%		340 20&
%		3\\
%		\hline
%		
%	\end{tabularx}
%\end{latin}}
% صورت اصلی مسئله را در این‌جا بنویسید.
%اگر مسئله چند قسمت داشت به این صورت عمل کنید:
%\شروع{زیرمسائل}
%\زیرمسئله[0] % مسئلۀ ۱آ
%\زیرمسئله % مسئلۀ ۱ب
%\زیرمسئله % مسئلۀ ۱پ
%\زیرمسئله % مسئلۀ ۱ت
%\پایان{زیرمسائل}




%%%%%%%%%%%%%%%%%%%%%%%%%%%%%%%%%%%%%%%%%%%%%%%%%%%%%
% برخی دستورات مفید برای تمرین‌ها. %
%%%%%%%%%%%%%%%%%%%%%%%%%%%%%%%%%%%%%%%%%%%%%%%%%%%%%

% فرمول‌های نمونه:
%$f(x) = \Theta(x)$
%$T(x, y) \leq \log(x) + 2^y + \binom{2n}{n}$
% {\tt code\_function}


% لیست موارد بدون شماره‌گذاری:
%\شروع{فقرات}
%    \فقره مورد اول 
%        \شروع{فقرات}
%            \فقره مورد اول داخلی 
%                \شروع{فقرات}
%                    \فقره اولین مورد خیلی داخلی 
%                    \فقره دومین مورد خیلی داخلی 
%                \پایان{فقرات}
%            \فقره مورد دوم داخلی 
%        \پایان{فقرات}
%    \فقره مورد دوم 
%\پایان{فقرات}

% لیست موارد شماره‌گذاری شده:
%\شروع{شمارش}
%    \فقره یک مورد 
%    \فقره یک مورد دیگر 
%    \فقره باز هم یک مورد دیگر
%\پایان{شمارش}

% برای نوشتن یک فرمول چند خطی که یک شماره بخوره. برای حالت‌های دیگر به سایت زیر مراجعه کنید:
% https://www.overleaf.com/learn/latex/Aligning_equations_with_amsmath
%\begin{align} 
%    \begin{split}
%        (x+y)^3 &= (x+y)^2(x+y) \\
%                &= (x^2+2xy+y^2)(x+y) \\
%                &= (x^3+2x^2y+xy^2) + (x^2y+2xy^2+y^3) \\
%                &= x^3+3x^2y+3xy^2+y^3
%    \end{split}                                 
%\end{align}

% نوشتن ماتریس و آرایه دوبعدی
%\begin{align}
%    A = 
%    \begin{bmatrix}
%        A_{11} & A_{21} \\
%        A_{21} & A_{22}
%    \end{bmatrix}
%\end{align}

% اضافه کردن تصویر با شرح. برای اضافه کردن تصویر بدون شرح، فقط دستور درج‌تصویر را استفاده کنید.
%\شروع{شکل}[h!]
%\تنظیم‌ازوسط
%  \درج‌تصویر[پهنا=0.5\پهنای‌متن]{img.jpg}
%\شرح{شرح تصویر}
%\برچسب{تصویر:۱}
%\پایان{شکل}
%\پایان{زیرمسائل}

%کد زیر را ببینید.
%%\lstinputlisting{source_filename.py}
%\begin{latin}
%\begin{lstlisting}
%# Import the modules
%import sys
%import random
%
%ans = True
%
%while ans:
%question = raw_input("Ask the magic 8 ball a question: (press enter to quit) ")
%
%answers = random.randint(1,8)
%
%if question == "":
%sys.exit()
%
%elif answers == 1:
%print "It is certain"
%
%elif answers == 2:
%print "Outlook good"
%
%elif answers == 3:
%print "You may rely on it"
%
%elif answers == 4:
%print "Ask again later"
%
%elif answers == 5:
%print "Concentrate and ask again"
%
%elif answers == 6:
%print "Reply hazy, try again"
%
%elif answers == 7:
%print "My reply is no"
%
%elif answers == 8:
%print "My sources say no"
%\end{lstlisting}
%\end{latin}
